
%----------------------------------------------------------------------------------------
%	Conclusion
%----------------------------------------------------------------------------------------
\section{Conclusion}
\subsection{Summary}

STARS is CubeSat intended to gather novel data about the inner Van Allen radiation belt.The on board payload is intended to gather data regarding the method by which high energy particles enter and exit the belt, their intensity, relative density, and location. The mission is split into three distinct phases, Pre-Mission, Mission, and Post-Mission. The Pre-Mission phase describes launch and placement into the desired orbit. The Mission phase encapsulates the time during which STARS will be gathering and transmitting scientific data. And the Post-Mission phase holds the end of life and mission termination procedures. A description of which of the individual subsystems will be active or inactive during each of these phases is provided.

Furthermore, STARS is broken up into six Critical Project Elements which define the most important aspects of the system. These six CPE's are Launch, Control, Power, Thermal, Communication and End of Life. These systems were deemed to be the most critical to the success of the mission and so it is these six CPE's upon which further trade studies and feasibility analysis will be performed.  Other aspects of the system, such as payload and mission analysis, will not have complete trade studies and are outlined in the Baseline Design Section.

12 functional requirements, defining the high level requirements for the STARS system, each with respective sub level deign requirements are also defined. The Key Design Alternatives for each CPE are then outlined and considered including a description of each of the possible options along with their benefits and drawbacks. This is followed by a description of the Trade Studies to be performed on each CPE. This includes a description of the different metrics by which each design alternative will be judged and of the weighting applied to each metric. An example of a possible Pugh matrix is also given as an indication of the results of such a trade study. This is followed by a description of the analysis to be performed on each CPE to ensure that the components selected via the trade study will perform as needed. The Baseline Design section includes a description of the parameters defines by the subsystems not includes in one of the six CPE's. This is followed by an outline of the project timeline which outlines the work that has already been completed along with the work to come. It includes an estimation of the time needed to accurately and correctly perform each task. This is followed with a diagram and discussion of the critical path in regard to further PDR level analysis. A brief discussion of budget considerations given the unique nature of this project is also included. 
\subsection{Discussions and Recommendations}

STARS is a complex system of smaller inter-reliant systems. The goal of STARS is to be able to attain new and interesting data regarding high energy particles in the Van Allen belts and then properly execute an end of life plan. In order to properly achieve its mission goals each of these subsystems must not only perform as designed but also properly interact with each of the other systems. As outlined in the sections above, significant analysis on all systems, particularly those classified as Critical Project Elements, needs to be performed to ensure their mutual functionality. This analysis will also need to be performed before down selection of specific individual pieces of each subsystem can be selected. A basic overview of the analysis intended to be performed is outlined in the Trade Study and Feasibility Study sections. The analysis is complex and often involves detailed simulations and feedback from various other systems. Once completed, NASC can begin the process of component selection and move on to more detailed and accurate analysis to be presented in the Critical Design Review document. 
